% TOP_PROBLEM: {"id": 2539, "title": "Kourovka Notebook Problem 21.30", "category_id": 17, "category": {"id": 17, "name": "group_theory", "display_name": "Group Theory", "description": "Problems about groups, group actions, representations, and related algebraic structures.", "slug": "group-theory", "order_index": 17, "created_at": "2026-07-31T15:26:25.671Z"}, "status": "open", "problem_number": "KOU-21.30", "set_id": 10}
% FIRST_POSTED: 2026-10-01
% ENTRY_KIND: statement_only
% UnsolvedMath ID 2539; source catalogue code KOU-21.30
% !TeX program = lualatex
% Definition and sources only; this file is not an LLM solution attempt.
\documentclass[11pt,a4paper]{article}
\usepackage[margin=25mm]{geometry}
\usepackage{fontspec}
\setmainfont{lmroman10-regular.otf}[BoldFont=lmroman10-bold.otf,
 ItalicFont=lmroman10-italic.otf,BoldItalicFont=lmroman10-bolditalic.otf]
\usepackage{amsmath,amssymb,mathtools}
\usepackage{unicode-math}
\setmathfont{latinmodern-math.otf}
\AtBeginDocument{\let\setminus\smallsetminus}
\usepackage{xurl,hyperref}
\hypersetup{colorlinks=true,urlcolor=blue}
\setcounter{secnumdepth}{0}
\setlength{\parindent}{0pt}
\setlength{\parskip}{5pt}
\setlength{\emergencystretch}{3em}
\begin{document}
\section*{Kourovka Notebook Problem 21.30}\label{2539}
{\small UnsolvedMath ID 2539 · KOU-21.30 · Group Theory · Kourovka Notebook - New Problems, Issue 21}
\subsection{Definitions and mathematical statement}
A finitely generated one-relator group is a group with a presentation
\[
 G=\langle x_1,\ldots,x_m\mid r=1\rangle
       =F_m/\langle\!\langle r\rangle\!\rangle,
\]
where $m\ge1$, $F_m$ is the free group on the displayed generators,
and $\langle\!\langle r\rangle\!\rangle$ is the smallest normal subgroup
containing the word $r$. The relator may be a proper power; no torsion-free
assumption is made. A redundant trivial relator causes no change to the
question.

A discrete group has the Haagerup property if it admits a metrically
proper action by affine isometries on a real Hilbert space $\mathcal H$,
possibly of infinite dimension. In detail, there are orthogonal operators
$\pi(g)$ and vectors $b(g)$ with
\[
 \pi(gh)=\pi(g)\pi(h),\qquad
 b(gh)=b(g)+\pi(g)b(h),
\]
so that $g\cdot v=\pi(g)v+b(g)$, and for every $R\ge0$ the set
$\{g\in G:\|b(g)\|\le R\}$ is finite. This is properness of the orbit
map of the origin; choosing another basepoint gives the same condition.

Does every one-relator group have this property? The Hilbert space and
action may depend on the group, and no uniform dimension or bound on
displacement is requested. Properness controls group elements, not merely
the set of orbit points. The affine formulation makes explicit the
source's isometric-action convention; group actions by bijective Hilbert
space isometries are affine in the real setting.
\subsection{Short English statement}
Does every group with a finite presentation containing one relator admit a proper affine isometric action on a Hilbert space?
\subsection{Sources}
\begin{itemize}
\item[S1] ULAM.AI and the UnsolvedMath Contributors, supplied problems.json, record 2539 (KOU-21.30), Kourovka Notebook - New Problems, Issue 21. Catalogue identity and source transcription; CC BY 4.0.
\url{https://huggingface.co/datasets/ulamai/UnsolvedMath}.
\item[S2] The Kourovka Notebook, 21st edition, new problems, Problem 21.30. Expanded formulation of the supplied catalogue statement.
\url{https://arxiv.org/pdf/1401.0300v47}.
\end{itemize}
\end{document}
